% Paper 3, §1 — Introduction: the reasoning triangle
% Anchors: kb/notes/reasoning/{chain-of-thought,test-time-compute,
%            inference-time-search,process-supervision,reasoning-training}.md
%          kb/notes/scaling/inference-time-compute-scaling.md
%          kb/notes/post-training/rlvr-and-grpo.md
%          kb/excerpts/{deepseek-r1,wei2022-cot,kojima2022,
%            wang2022-self-consistency,snell2024,lightman2023-prm800k,
%            shao2024,dapo2025,s1-2025,rstar-math2025}.md
% Cross-paper: Paper 1 §architecture (long-context attention, KV-cache cost
%              of long traces); Paper 2 §RL-stage (RLHF→RLVR transition).

\section{Introduction: the reasoning triangle}
\label{sec:introduction}

Between the late-2022 release of \citet{wei2022-cot} and the early-2025
release of \citet{deepseek-r1}, the publicly reported pass@1 of an
open-weight model on AIME 2024 --- a contest of fifteen graduate-style
mathematics problems --- moved from indistinguishable-from-noise to
$77.9\%$, and to $86.7\%$ under majority vote over $64$ samples
\citep[\S 2.3]{deepseek-r1}.\footnote{KB excerpt:
\texttt{kb/excerpts/deepseek-r1.md\#sec-2-3-aime}.} The same checkpoint,
DeepSeek-\glsterm{r1-zero}{R1-Zero}, started its reinforcement-learning phase from a
$15.6\%$ baseline on the same benchmark
\citep[\S 2.3]{deepseek-r1}.\footnote{KB excerpt:
\texttt{kb/excerpts/deepseek-r1.md\#sec-2-3-aime}.} The size of the
jump --- and the comparable jumps reported by OpenAI's o1, Google's
Gemini~2.5-thinking, Anthropic's \glsterm{claude}{Claude}~3.7 extended-thinking mode,
Alibaba's QwQ, and the Qwen3-thinking series in the months following
--- is the clearest empirical event in language-model research between
2022 and 2026. It is also the event around which the literature is
most prone to one-leg explanations: ``the model thinks in chains''; ``a
verifier ranks its samples''; ``a reinforcement-learning loop trains
it.'' Each one-leg story is partially right and structurally
incomplete.

The thesis of this paper is that the 2022--2026 reasoning gains are
best understood as a \emph{triangle of co-designed levers}. The three
legs:

\begin{enumerate}
\item \textbf{Inference-time compute} — \glsterm{chain-of-thought}{chain-of-thought}
  \citep{wei2022-cot,kojima2022}, sampling-based aggregation
  \citep{wang2022-self-consistency}, and the elevation of test-time
  compute from prompting trick to scaling axis by
  \citet{snell2024}.\footnote{KB excerpt:
  \texttt{kb/excerpts/snell2024.md\#sec-1-flops-matched}.}

\item \textbf{Verifier signals} — process reward models that score
  partial reasoning at step granularity \citep{lightman2023-prm800k}
  and the tree-search families (beam, BFS, MCTS) that turn those
  scores into search trajectories \citep{rstar-math2025}.

\item \textbf{Reinforcement learning on verifiable rewards (\glsterm{rlvr-3}{RLVR})} ---
  group-relative policy optimisation over programmatically-checkable
  outcomes \citep{shao2024}, with \glsterm{dapo-2}{DAPO}'s four targeted modifications
  \citep{dapo2025} as the current open-recipe state of the art.
\end{enumerate}

Frontier reasoning models in 2026 ship all three. They differ in
mixture, in budget, in disclosure, and in which leg is treated as the
primary product surface (R1 ships the long-\glsterm{chain-of-thought}{CoT} trace; o1 ships a
hidden reasoning channel and an exposed effort knob; Gemini 2.5
exposes a literal \texttt{thinkingBudget}). They do not, as of
2026-Q2, differ in shipping fewer than three legs. Treating any one
leg in isolation reproduces an incomplete reading of the headline
gains, of which mechanism the gains generalise from, and of which open
questions remain.

\begin{intuition}
A useful one-line picture is that the three legs are dual to each
other across the train/inference boundary. \glsterm{rlvr-3}{RLVR} \emph{shapes} a policy
whose long-\glsterm{chain-of-thought}{CoT} output is worth amplifying; inference-time compute
\emph{amplifies} that policy by spending more sampling, search, or
trace-length budget; verifier signals \emph{rerank} the amplified
trajectories. Pull any one leg and the other two lose a lot of their
explanatory power: long-\glsterm{chain-of-thought}{CoT} without an RL phase that taught the model
to use it is short-\glsterm{chain-of-thought}{CoT} in a costume; a strong verifier without a
proposer that produces diverse-and-mostly-right candidates does not
have anything to verify; an RL loop without a \glsterm{verifiable-reward}{verifiable reward}
collapses back to \glsterm{rlhf}{RLHF} and its 2023-vintage pathologies. Returning to
canonical math: the \glsterm{grpo-2}{GRPO} objective transcribed in
\cref{sec:rlvr-grpo-objective} is structurally a re-weighted
maximum-likelihood update over verifier-passed trajectories sampled
from the current policy, with KL anchoring to a reference. All three
legs appear in that single objective.
\end{intuition}

The triangle framing also locates this paper relative to two of its
companions. Paper~2 establishes a six-stage post-training pipeline in
which \glsterm{rlhf}{RLHF} is succeeded by RL-on-verifiable-rewards as a distinct
stage, with different load-bearing variables (programmatic verifier
instead of learned reward, group baseline instead of \glsterm{value}{value} network).
That cross-reference is the inheritance route: \cref{sec:rlvr-grpo-dapo}
deepens Paper~2's \glsterm{rlvr-3}{RLVR} treatment in the reasoning-specific direction
(long traces, terminal verifiers, R1-style protocols). Paper~1 is the
architectural precondition: the 10K-token reasoning traces that
\cref{sec:r1-family} reports as the operational signature of an
RL-trained \glsterm{reasoning-model}{reasoning model} are tractable only because Paper~1's
long-context attention (sliding-window, \glsterm{flashattention}{FlashAttention}-2/3,
KV-compression) makes their \glsterm{kv-cache}{KV cache} survivable.\footnote{The
KV-cache cost of a $10^4$-token trace at full $\BHSdh$ shape is
non-trivial. The architectural choices of Paper~1 are what determine
whether a frontier \glsterm{reasoning-model}{reasoning model} is feasible as a deployment
artefact, not only as a research curiosity.} Without those two
predecessors the triangle's three legs are a \glsterm{vocabulary}{vocabulary}; with them,
they are a co-designed system.

\subsection{What this paper does, section by section}
\label{sec:intro-roadmap}

The body of the paper walks the triangle from each leg's historical
entry point through its anchor mechanisms to the live contradictions
the literature has not yet closed.

\Cref{sec:cot-phenomenon} starts with the phenomenon: \citet{wei2022-cot}
and \citet{kojima2022} on \glsterm{chain-of-thought}{chain-of-thought} as a few-shot and zero-shot
prompting recipe, the decoding factorisation $z \sim p_\theta(\cdot
\mid x), y \sim p_\theta(\cdot \mid x, z)$ that every later
test-time-compute strategy in the paper inherits, and the
emergence-with-scale finding that distinguished \glsterm{chain-of-thought}{CoT} from
compute-monotonic prompting tricks. \Cref{sec:self-consistency}
continues with the cheapest extension: \citet{wang2022-self-consistency}'s
sample-$N$-and-majority-vote, no verifier needed --- still a competitive
baseline well into the R1 era, where it lifts \glsterm{r1-zero}{R1-Zero}'s AIME 2024 from
$77.9\%$ to $86.7\%$ \citep[\S 2.3]{deepseek-r1}.

\Cref{sec:inference-compute-scaling} is the first anchor section: the
\citet{snell2024} reframing of \glsterm{test-time-compute}{test-time compute} as a scaling axis with
its own functional form, comparable to Chinchilla on the training
side. The compute-optimal definition (Eq.~1 of
\citet{snell2024}), the proposer/verifier decomposition that unifies
every other section of the paper, and the FLOPs-matched headline that
\glsterm{test-time-compute}{test-time compute} substitutes for $\sim\!14\times$ pretraining FLOPs on
easy and intermediate prompts but not on the hardest tail, all live
there.

\Cref{sec:process-reward-models} formalises the verifier leg.
\citet{lightman2023-prm800k}'s process \glsterm{reward-model}{reward model} is the canonical
artefact; the four trace-aggregation rules (min, product, last-step,
weighted-mean) and the headline that min-aggregation is strongest are
the load-bearing mechanics. The post-2023 \glsterm{process-reward-model}{PRM} lineage --- Math-Shepherd
on automated step labelling, \glsterm{generative-prm}{ThinkPRM} and R-\glsterm{process-reward-model}{PRM} on generative
verifiers that themselves reason --- closes the section.
\Cref{sec:tree-search} is the second anchor section: the search-tree
datastructure over partial reasoning prefixes, the canonical \glsterm{reasoning-mcts}{reasoning-MCTS}
recurrence transcribed from \citet{rstar-math2025}, the four MCTS
phases (selection, expansion, simulation, backup), and ReST-MCTS\textsuperscript{*}
as the AlphaZero-style outer loop in which MCTS rollouts harvest
correct trajectories that re-train the proposer and the verifier.

\Cref{sec:rlvr-grpo-dapo} is the third anchor section and the third
triangle leg. The verifier reward
$R_\mathrm{verifier}(x,y) = \mathbb{1}[\,\mathrm{verify}(x,y) =
\mathsf{true}\,]$, the \glsterm{grpo-2}{GRPO} objective \citep[\S 4.1]{shao2024} with
group-relative advantage replacing a \glsterm{value}{value} network, and \glsterm{dapo-2}{DAPO}'s four
targeted modifications \citep[\S 2]{dapo2025} are the section's
mechanics. The empirical proof point that \cref{sec:r1-family} unfolds
in detail --- DeepSeek-\glsterm{r1-zero}{R1-Zero}'s continuous \glsterm{grpo-2}{GRPO} climb --- is the
cleanest single number in the literature for the proposition that the
three legs compose under one training loop.

\Cref{sec:budget-forcing} treats test-time scaling that does not lean
on parallel sampling or search: \citet{s1-2025}'s minimum recipe of
$1{,}000$ long-\glsterm{chain-of-thought}{CoT} exemplars plus a budget-forcing rule that appends
\texttt{Wait} to extend thinking, isolating the
\emph{reasoning-architecture} axis (the \texttt{<think>...</think>}
phase as a structurally distinct decode-time region) from the
base-architecture axis. \Cref{sec:search-vs-rl} confronts the
operational tradeoff the triangle implies: at fixed total compute,
when does the marginal FLOP go to inference search, to \glsterm{rlvr-3}{RLVR} training,
or to pretraining? The honest answer in 2026 is benchmark-specific and
base-model-specific; that section transcribes what is settled and
demarcates what is anecdotal.

\Cref{sec:cot-faithfulness} addresses the most consequential live open
question: does the trace $z$ causally drive the answer $y$, or is it
post-hoc rationalisation that happens to be human-readable? The
2025--2026 wave (Arcuschin et al.\ on naturalistic prompts, Shen et
al.\ on FaithCoT-Bench, Mehta et al.\ on whether faithfulness improves
with scale) finds substantial unfaithfulness on naturalistic
distributions; whether it improves with scale is contested.
\Cref{sec:prm-survival} returns to the verifier leg with a sharper
question: post-R1, do PRMs still pay for themselves? R1's headline
result was achieved with terminal verifier rewards only, and the
2025--2026 \glsterm{process-reward-model}{PRM} literature is split between diminishing-returns
findings and continued-gain findings.

\Cref{sec:r1-family} sits between \cref{sec:rlvr-grpo-dapo} and
\cref{sec:budget-forcing} as the empirical proof point: the four-stage
R1 protocol (cold-start SFT $\to$ \glsterm{grpo-2}{GRPO} with verifiable rewards $\to$
\glsterm{rejection-sampling-sft}{rejection-sampling SFT} on harvested traces $\to$ general-domain RL),
the \glsterm{r1-zero}{R1-Zero} ablation that strips the cold-start to isolate what the
SFT initialisation contributes, and the open-weight distillation track
that produced the post-R1 reasoning ecosystem (R1-Distill-Qwen-7B/14B/32B,
the Tülu~3 lineage, OpenR1, QwQ, Qwen3-thinking).

The final two sections close the loop. \Cref{sec:frontier-snapshot}
catalogues which 2026-frontier models ship which combination of legs
at what budget; the table is intended as the most-screenshotted figure
of the paper. \Cref{sec:contradictions} concentrates the
\textsf{[CONTRADICTION]} markers from \cref{sec:cot-phenomenon}--\cref{sec:prm-survival}
into one place and forecasts which of them are likely to close in the
next model generation.

\subsection{The open question this paper forwards}
\label{sec:intro-open-question}

The thesis commits the paper to one position --- the legs are
co-designed in the strong sense, not three independent gains that
happen to ship together --- and one open question that
\cref{sec:contradictions} returns to:

\begin{quote}
\emph{Is the triangle a real co-design, or three independent gains we
happen to ship together?}
\end{quote}

The strong-co-design claim is that the legs interlock causally: \glsterm{rlvr-3}{RLVR}
without long-\glsterm{chain-of-thought}{CoT} sampling cannot find the signal that lengthens
traces, long-\glsterm{chain-of-thought}{CoT} without \glsterm{rlvr-3}{RLVR} is mechanistically possible but
empirically weaker, and verifier signals are what convert sampling
diversity into a reranked answer that propagates training reward.
The independent-gains claim is that the three legs each work in
isolation and the field has co-shipped them only for engineering
reasons. The 2025 \citet{rlvr-limits-2025} result --- \glsterm{rlvr-3}{RLVR} sharpens
pass@1 within base-model competency but does not expand pass@$K$ ---
sits exactly on this fault line: under the strong-co-design reading,
\glsterm{rlvr-3}{RLVR} is doing the right thing (it shapes a policy whose sampling and
verification are then meaningful); under the independent-gains
reading, \glsterm{rlvr-3}{RLVR} is doing less than the headline numbers suggest because
the underlying capability already exists in the base model and only
the trace-length needed to express it is being elicited.
\Cref{sec:contradictions} catalogues five further contradictions on
the same axis (\glsterm{process-reward-model}{PRM} utility post-R1, \glsterm{cot-faithfulness-2}{CoT faithfulness} scaling,
search-vs-RL tradeoff geometry, AIME contamination, MCTS-vs-beam at
matched budget). The paper's position is that none of them are closed
as of 2026-Q2; the engineering question of how to ship a frontier
\glsterm{reasoning-model}{reasoning model} is, by contrast, well past the point where any single
leg is optional.

The remainder of the paper makes that position defensible at
mechanism, equation, and empirical-benchmark granularity. We begin in
\cref{sec:cot-phenomenon} with the historical entry point that every
later leg of the triangle inherits: the \glsterm{chain-of-thought}{chain-of-thought} decoding
pattern.
